% SGA 6，PDF 1--29：书前材料与第0讲；不含第0讲附录。
\SourcePage{1}
\begin{center}
{\Large 数学讲义丛书}\par
非正式报告与讨论班文集\par
主编：A. Dold（海德堡）、B. Eckmann（苏黎世）\par
第225卷\par\bigskip
玛丽森林\Term{algebraic-geometry}{代数几何}讨论班，1966/67\par
SGA 6\par
主持：P. Berthelot、A. Grothendieck、L. Illusie\par
参与者：D. Ferrand、J. P. Jouanolou、O. Jussila、S. Kleiman、M. Raynaud、J. P. Serre\par\bigskip
{\Large \Term{intersection-theory}{相交理论}与黎曼--罗赫\Term{theorem}{定理}}\par
Springer-Verlag，柏林、海德堡、纽约
\end{center}
\SourcePage{2}
\begin{center}
数学讲义丛书\par
非正式报告与讨论班文集\par
主编：A. Dold（海德堡）、B. Eckmann（苏黎世）\par
第225卷\par\bigskip
\Term{intersection-theory}{相交理论}与黎曼--罗赫\Term{theorem}{定理}\par
Springer-Verlag\par
柏林、海德堡、纽约，1971年
\end{center}
\SourcePage{3}
AMS主题分类（1970）：14-02，14C15，14F15，14K30，18E30。

ISBN 3-540-05647-5，Springer-Verlag，柏林、海德堡、纽约。\par
ISBN 0-387-05647-5，Springer-Verlag，纽约、海德堡、柏林。

本书受版权保护。无论涉及全部内容还是部分内容，一切权利均予保留，尤其包括翻译、重印、插图再利用、广播、以复印机或类似方法复制，以及存入数据库的权利。根据德国版权法第54条，出于私人使用以外的目的制作复制品时，须向出版社支付费用，金额由双方协商确定。

\copyright\ Springer-Verlag，柏林、海德堡，1971年。美国国会图书馆目录卡号73-178287。德国印制。胶版印刷：Julius Beltz，Hemsbach/Bergstr。
\SourcePage{4}
\begin{center}
玛丽森林\Term{algebraic-geometry}{代数几何}讨论班\par
1966--67\par\bigskip
{\Large \Term{intersection-theory}{相交理论}与黎曼--罗赫\Term{theorem}{定理}}\par
（SGA 6）\par\bigskip
讨论班主持\par
P. Berthelot、A. Grothendieck、L. Illusie\par\bigskip
参与者\par
D. Ferrand、J. P. Jouanolou、O. Jussila、S. Kleiman、M. Raynaud、J. P. Serre
\end{center}
\SourcePage{5}
\SGAsection{}{序言}
本版《\Term{algebraic-geometry}{代数几何}讨论班6》与最初由IHES以油印分讲\Term{form}{形式}发行的版本基本相同。所作改正主要针对印刷错误；但第I讲和第III讲例外，L. Illusie对初版中若干不正确的陈述作了轻微修改。此外，还添加了一些脚注，尤其是A. Grothendieck在第XIV讲中添加的脚注，以说明讨论班举行时尚未解决的问题的最新进展。最后还须指出，D. Ferrand关于\Term{perfect-complex}{完美复形}的\Term{determinant-functor}{行列式}理论的第XI讲并未写成，因而不会出版；为免引用发生错误，其后各讲的编号仍原样保留。
\begin{flushright}1971年6月\end{flushright}
\SourcePage{6}
\SGAsection{}{引言}
本讨论班的目的是为任意\Term{scheme}{概形}建立整体的\Term{intersection-theory}{相交理论}和一个黎曼--罗赫公式。第0讲介绍本讨论班的纲领。沿着Borel--Serre于1958年发表的著名报告的思路，为\Term{generic}{一般}\Term{regular-scheme}{正则概形}证明一个黎曼--罗赫公式，在原则上从那时起就已经是清楚的，至少对\Term{projective-morphism}{射影态射}如此。向\Term{generic}{一般}情形的推广则不那么显然；包含这项推广的纲领（本讨论班部分地实现了它）可以追溯到1960年。如同\Term{coherent-sheaf}{凝聚层}的\Term{duality-theory}{对偶理论}一样（参见R. Hartshorne，\textit{Residues and Duality}，数学讲义丛书第20卷，Springer），要建立一个令人满意的理论，Verdier的\Term{derived-category}{导出范畴}理论是必不可少的工具；了解这一理论，是读懂本讨论班所必需的。除此之外，研读本讨论班几乎只需对EGA第I、II、III章所述的\Term{algebraic-geometry}{代数几何}基础有\Term{generic}{一般}了解。我们还将偶尔引用EGA IV中关于\Term{flat-morphism}{平坦}或\Term{smooth-morphism}{光滑态射}的若干事实，而这些事实的主要部分也已出现在SGA 1的前几讲。最后，为了使某些结果达到应用所需的\Term{generic}{一般}程度，我们有时使用\Term{ringed-site}{赋环景}和\Term{ringed-topos}{赋环拓扑斯}的概念，相关内容参见SGA 4第I至IV讲。不熟悉景与\Term{topos}{拓扑斯}语言的读者，可以处处将这些\Term{object}{对象}换成通常的\Term{topological-space}{拓扑空间}，将\Term{topos}{拓扑斯}的\Term{object}{对象}换成这些\Term{space}{空间}的\Term{open-subset}{开集}；不过我们仍更建议读者掌握\Term{topos}{拓扑斯}的语言，它提供了极其方便的统一原则。

这里所述理论的基本概念是\Term{perfect-complex}{模的完美复形}及其各种变体；这些概念将在第I、II、III讲中展开。显然，它们在\Term{algebraic-geometry}{代数几何}的其他场合也很重要（尤其是\Term{discrete}{离散}系数\Term{cohomology}{上同调}〔SGA 5〕或\Term{coherent}{凝聚}系数\Term{cohomology}{上同调}中的Lefschetz--Verdier型公式），而且还将在\Term{algebraic-geometry}{代数几何}之外发挥\Term{action}{作用}，特别是在表述这里的黎曼--罗赫\Term{theorem}{定理}的复解析变体，或Atiyah--Singer指标\Term{theorem}{定理}的适当变体时。
\SourcePage{7}
第II讲将给出这方面的一些提示。遗憾的是，目前尚无一个陈述，哪怕是启发性的陈述，能够涵盖这两个\Term{theorem}{定理}（其中第一个目前仍是猜想）。看来这里还缺少一个新思想；同样，在\Term{algebraic-geometry}{代数几何}中，若要在没有\Term{projective}{射影}性假设的情形下证明黎曼--罗赫\Term{theorem}{定理}，也缺少这样的新思想（参见第XIV讲第2节）。

本讨论班完全不会使用J. P. Serre的《\Term{locally-algebraic}{局部代数}·\Term{multiplicity}{重数}》（数学讲义丛书第11卷，Springer）所述的\Term{local-scheme}{局部}\Term{intersection-theory}{相交理论}；这里只指出，Serre的这门课程显然影响了本理论于1957年的产生。我们也不使用1958年Chevalley讨论班第2、3讲（巴黎高等师范学校）所建立的\Term{chow-ring}{Chow环}理论。该理论本质上依赖于\Term{nonsingular}{非奇异}性假设，而本讨论班的目标恰是为\Term{generic}{一般}\Term{scheme}{概形}（乃至\Term{generic}{一般}\Term{locally-ringed-topos}{局部赋环拓扑斯}）建立\Term{intersection-theory}{相交理论}；关于这一点，参见第XIV讲第4、8节的评述，那里说明了我们的理论与Chow环理论之间的关系，并提出了推广后者的问题。可以说，本讨论班给出了一个相容而自足的\Term{intersection-theory}{相交理论}，但它绝未穷尽\Term{intersection-theory}{相交理论}中有用（甚至不可缺少）的各种观点，因此适宜以所引Serre与Chevalley的讲稿作为补充。

我们把Grothendieck于1957年写成的报告《层的类与黎曼--罗赫\Term{theorem}{定理}》附在本讨论班之后（作为第0讲的附录）。该报告曾是同年Borel与Serre在普林斯顿举办讨论班，以及他们上述论文的基础。报告概述了\Term{nonsingular}{非奇异}\Term{quasi-projective-morphism}{拟射影}簇的黎曼--罗赫\Term{theorem}{定理}的两种证明；其中第一种目前只适用于\Term{characteristic-zero}{特征零}，但对于\Term{immersion}{浸入}给出的结果稍为精确，这种证明尚未见诸文献（Borel--Serre的工作和本讨论班也包括在内）。当然，阅读这份报告不以读过本讨论班的任何其他一讲为前提；它甚至可以与第0讲一样，作为学习本讨论班的入门，尤其适合尚不熟悉Borel--Serre工作的读者。刚才提到的证明，本质上可以原封不动地用于\Term{complex-analytic-space}{复解析空间}的\Term{projective-morphism}{射影态射}；
\SourcePage{8}
一旦\Term{derived-category}{导出范畴}中的“\Term{exterior-power}{外幂}”运算问题得到解决（参见第XIV讲第1、2、3节），它无疑也将适用于任意\Term{characteristic}{特征}。即使从本讨论班所采取的观点来看，没有研究这个问题，也无疑是本讨论班最严重的缺憾。

读者会注意到，扉页所列的两位参加者没有出现在本讨论班的目录中。一位是J. P. Jouanolou，他积极参与了讨论班第一部分的技术准备工作，但未能参加口头报告；重要陈述VI 1.1中关于\Term{projective-bundle}{射影丛}上\Term{locally-free-module}{局部自由模}的类环$K'$的结构所包含的\Term{linearly-independent}{线性无关}性断言，尤其应归功于他（此前只有在\Term{base-scheme}{底概形}有一个\Term{ample-invertible-module}{丰沛可逆模}的假设下才证明过）。另一位是J. P. Serre，他作过两次口头报告，内容在逻辑上独立于本讨论班其余部分，因此他选择将它们另行发表为论文（参见J. P. Serre，\textit{Groupes de Grothendieck des schémas en groupes réductifs déployés}，拟刊于\textit{Publications Mathématiques}第34号）。

如同玛丽森林\Term{algebraic-geometry}{代数几何}讨论班的其他部分，简称SGA 1至SGA 6指该讨论班的不同部分，紧随其后的罗马数字为讲次，后面的阿拉伯数字则对应于该讲内部的编号。对于本讨论班内部的引用，省略简称SGA 6，使用I 4.9这样的写法（意为第I讲，陈述4.9）。简称EGA指J. Dieudonné与A. Grothendieck的《\Term{algebraic-geometry}{代数几何}基础》。
\SourcePage{9}
\SGAsection{}{目录}
% 原目录PDF9--10；页码由译本总目录重建。以下保留完整题名、作者与原书起页信息。
\begin{description}
\item[第0讲] \Term{generic}{一般}\Term{scheme}{概形}上\Term{intersection-theory}{相交理论}的纲领概述，A. Grothendieck（原书第1页）。
\item[附录] 层的类与黎曼--罗赫\Term{theorem}{定理}，A. Grothendieck（原书第20页）。
\item[第I讲] \Term{derived-category}{导出范畴}中\Term{finiteness-condition}{有限性条件}的\Term{generic}{一般}理论，L. Illusie（原书第78页）。
\item[第II讲] 整体\Term{resolution}{消解}的存在性，L. Illusie（原书第160页）。
\item[第III讲] 相对\Term{finiteness-condition}{有限性条件}，L. Illusie（原书第222页）。
\item[第IV讲] \Term{ringed-topos}{赋环拓扑斯}的Grothendieck群，L. Illusie（原书第274页）。
\item[第V讲] $\lambda$环的\Term{generic}{一般}理论，P. Berthelot（原书第297页）。
\item[第VI讲] \Term{projective-bundle}{射影丛}的$K'$：计算与推论，P. Berthelot（原书第365页）。
\end{description}
\SourcePage{10}
\begin{description}
\item[第VII讲] \Term{regularly-immersed}{正则浸入}与\Term{blow-up}{爆破}\Term{scheme}{概形}的$K'$的计算，P. Berthelot（原书第416页）。
\item[第VIII讲] 黎曼--罗赫\Term{theorem}{定理}，P. Berthelot（原书第466页）。
\item[第IX讲] $K_{\bullet}$群的若干计算，P. Berthelot（原书第498页）。
\item[第X讲] \Term{proper}{固有}\Term{algebraic-scheme}{代数概形}上的相\Term{intersection-form}{交形式}理论，O. Jussila；附A. Grothendieck所作附录《\Term{intersection-theory}{相交理论}中的\Term{specialization}{特殊化}》（原书第519页）。
\item[第XI讲] 未写成。
\item[第XII讲] 关于Picard\Term{functor}{函子}的一个相对\Term{representability}{可表性}\Term{theorem}{定理}，M. Raynaud（S. Kleiman整理）（原书第595页）。
\item[第XIII讲] Picard\Term{functor}{函子}的\Term{finiteness-theorem}{有限性定理}，S. Kleiman（原书第616页）。
\item[第XIV讲] \Term{intersection-theory}{相交理论}中的未解决问题，A. Grothendieck（原书第667页）。
\item[术语索引] 原书第691页。
\item[记号索引] 原书第696页。
\end{description}
\SourcePage{11}
\SGAchapter{0}{一般概形上相交理论的纲领概述}{Esquisse d'un Programme pour une Théorie des Intersections sur les Schémas Généraux}
\begin{center}A. Grothendieck\end{center}
本讲具有导论性质；在逻辑上，学习本讨论班并不必先读本讲。它尤其面向熟悉黎曼--罗赫\Term{theorem}{定理}的暂行版本的读者，该版本见Borel--Serre的报告[2]，或引言中已提到的Grothendieck的报告（以下记作[RRR]），后者作为附录重印于本讲之后。
\SGAsection{1}{}
先回忆如下黎曼--罗赫公式。设
\[f:X\longrightarrow Y\]
是域$k$上\Term{quasi-projective-scheme}{拟射影}且\Term{smooth-morphism}{光滑的概形}之间的\Term{proper-morphism}{固有态射}，$F$是$X$上的\Term{coherent-sheaf}{凝聚层}，其在$X$上\Term{coherent-sheaf}{凝聚层}的类群$K(X)$中的类记为$\operatorname{cl}(F)$。则
\[
\operatorname{Todd}(T_Y)\operatorname{ch}_Y\bigl(f_*(\operatorname{cl}(F))\bigr)
=f_*\bigl(\operatorname{Todd}(T_X)\operatorname{ch}_X(F)\bigr).
\tag{1.1}
\]
此公式在$A(Y)\otimes_{\mathbb Z}\mathbb Q$中成立，其中$A(Y)$是$Y$的\Term{chow-ring}{Chow环}。右边的$f_*$由\Term{group-homomorphism}{群同态}“\Term{direct-image}{圈的直像}”
\[f_*:A(X)\longrightarrow A(Y)\]
张量以$\mathbb Q$得到；左边的$f_*$则是$F$关于$f$的\Term{euler-characteristic}{欧拉--庞加莱示性数}：
\[f_*(\operatorname{cl}(F))=\sum_i(-1)^i\operatorname{cl}(R^if_*(F)).\]
$\operatorname{ch}_X$、$\operatorname{ch}_Y$分别表示$X$、$Y$上的\Term{chern-character}{陈特征}，$T_X$、$T_Y$分别为$X$、$Y$的\Term{tangent-bundle}{切丛}。众所周知，$\operatorname{Todd}(-)$与$\operatorname{ch}(-)$是以其自变量的\Term{chern-class}{陈类}为变量、有理系数的\Term{universal-polynomial}{泛多项式}。$\operatorname{Todd}(-)$的常数项为1，故对任意自变量值它都是\Term{invertible-element}{可逆元}。因此，将(1.1)乘以$\operatorname{Todd}(T_Y)^{-1}$，可得更适合于我们目的的\Term{equivalence}{等价}\Term{form}{形式}
\[
\operatorname{ch}_Y\bigl(f_*(\operatorname{cl}(F))\bigr)
=f_*\bigl(\operatorname{Todd}(T_f)\cdot\operatorname{ch}_X(F)\bigr),
\tag{1.2}
\]
其中置
\[T_f=T_X-f^*(T_Y)\in K(X),\tag{1.3}\]
\SourcePage{12}
从而$T_f$起着$X$在$Y$上的\Term{virtual-relative-tangent-bundle}{虚相对切丛}的\Term{action}{作用}。当\Term{morphism}{态射}$f$\Term{smooth}{光滑}时（即\Term{tangent-map}{切映射}处处\Term{surjectivity}{满射}），直接有
\[T_f=T_{X/Y}\qquad\text{（沿纤维方向的切丛）},\]
而当$f:X\longrightarrow Y$是\Term{immersion}{浸入}时，则得到
\[T_f=-\check N_{X/Y},\]
其中$\check N_{X/Y}$表示$X$在$Y$中的\Term{normal-module}{法层}。

本讨论班的主要目标之一，是同时从两个方向推广(1.2)：
\begin{enumerate}
\item[a)] 去掉存在底域$k$的假设。
\item[b)] 将$Y$、$X$的\Term{regularity}{正则性}假设换成$f$的某种“\Term{local-scheme}{局部}\Term{regularity}{正则性}”假设。
\end{enumerate}
最后，在此过程中，我们还将处理如下问题：
\begin{enumerate}
\item[c)] 消去\Term{quasi-projective-morphism}{拟射影}性假设；在没有底域时，这一假设表述为$X$、$Y$上存在\Term{ample-invertible-module}{丰沛可逆模}。
\end{enumerate}
\SGAsection{2}{}
先考察推广a)，但仍保留$X$、$Y$\Term{regular}{正则}且具有\Term{ample-invertible-module}{丰沛可逆模}的假设。

由于$X$、$Y$\Term{regular}{正则}，$K(X)$、$K(Y)$以及\Term{homomorphism}{同态}$f_*:K(X)\longrightarrow K(Y)$的定义没有带来新的困难。因此，赋予(1.2)意义的最自然途径似乎是：定义\Term{chow-ring}{Chow环}$A(X)$、$A(Y)$及\Term{group-homomorphism}{群同态}
\[f_*:A(X)\longrightarrow A(Y);\]
建立相应的\Term{chern-class}{陈类}理论，给出映射
\[c_i:K(X)\longrightarrow A(X),\]
以及$Y$上的相应映射；最后，再描述一个\Term{virtual-relative-tangent-bundle}{虚相对切丛}元素
\[T_f\in K(X).\]
\SourcePage{13}
\textbf{2.1.} 关于最后这一项，一个定义相当明显地呈现出来。利用$X$、$Y$上存在\Term{ample-invertible-module}{丰沛可逆模}这一假设，可以将\Term{morphism}{态射}$f$\Term{resolution}{分解}为
\[X\xrightarrow{i}X'\xrightarrow{f'}Y,\tag{2.1}\]
其中$i$是\Term{closed-immersion}{闭浸入}，$f'$是\Term{smooth-morphism}{光滑态射}；例如，对适当的$r$可取$X'=\mathbb P^r_Y$。于是置
\[\check T_f=\operatorname{cl}(\Omega^1_{X'/Y})-\operatorname{cl}(N_{X/X'}),\tag{2.2}\]
其中$\Omega^1_{X'/Y}$为$X'$在$Y$上的相对一次\Term{differential}{微分}的\Term{locally-free-module}{局部自由模}（或称$X'$在$Y$上的相对余切模），$N_{X/X'}$为$X$在$X'$中的\Term{conormal-module}{余法模}，按定义等于$J/J^2$；$J$是$X'$上定义\Term{closed-subscheme}{闭子概形}$X$的\Term{ideal}{理想}。由于$X$、$X'$\Term{regular}{正则}，还可\Term{pushout}{推出}$N_{X/X'}$也是\Term{locally-free}{局部自由}的。最后，$\check T_f$表示$T_f$的\Term{duality}{对偶}。不难验证（第VIII讲），由(2.1)定义的元素$T_f$与所选的$f$的\Term{resolution}{分解}无关。

\textbf{2.2.} 至于在上述条件下定义\Term{chow-ring}{Chow环}及相应\Term{chern-class}{陈类}理论，虽然看来应当能仿照\Term{smooth}{光滑}\Term{quasi-projective-morphism}{拟射影}\Term{algebraic-scheme}{代数概形}的已知理论来建立，但目前似乎尚未完成这项工作，本讨论班当然也未着手于此。相反，我们引入另一个环，它将起着类似于\Term{chow-ring}{Chow环}的\Term{action}{作用}；当$X$在域$k$上\Term{quasi-projective-morphism}{拟射影}且\Term{smooth}{光滑}时，这个环与$\mathbb Q$的\Term{tensor-product}{张量积}确实\Term{isomorphism}{同构}于$A(X)\otimes_{\mathbb Z}\mathbb Q$。这个环是对$K(X)$施加适当\Term{filtered}{滤过}后所得的\Term{associated-graded-ring}{相伴分次环}。

$K(X)$上第一个自然的\Term{filtration}{滤过}，即[RRR]所考察的\Term{filtered}{滤过}，是按层的\Term{support}{支集}的\Term{codimension}{余维数}给出的“\Term{topological-filtration}{拓扑滤过}”：$\operatorname{Filt}^i_{\mathrm{top}}K(X)$由所有$\operatorname{cl}(F)$生成，其中$F$\Term{coherent}{凝聚}，且$\operatorname{codim}(\operatorname{supp}F,X)\geq i$。第一个问题是，这个\Term{filtered}{滤过}是否与\Term{ring-structure}{环结构}相容；当$X$在域$k$上\Term{smooth}{光滑}时，[RRR]借助Chow的“\Term{moving-lemma}{移动引理}”解决了这个问题。承认这一点后，就得到一个\Term{associated-graded-ring}{相伴分次环}，这里记作$\operatorname{Gr}'_{\mathrm{top}}(X)$。[RRR]定义了一个满的\Term{ring-homomorphism}{环同态}
\[\phi:A(X)\longrightarrow\operatorname{Gr}'_{\mathrm{top}}(X),\]
\SourcePage{14}
它将素圈$Z$的类变为$\operatorname{cl}(\mathcal O_Z)$；后来发现（第XIV讲4.2），其核为\Term{torsion-subgroup}{挠子群}。此外，[RRR]说明，\Term{chern-class}{陈类}$c_i(F)$在$\operatorname{Gr}'_{\mathrm{top}}(X)$中的像，可以直接用$\operatorname{cl}(F)\in K(X)$和$K(X)$中的\Term{exterior-power}{外幂}运算$\lambda^i$计算。

这说明，对于具有\Term{ample-invertible-module}{丰沛可逆模}的\Term{noetherian}{诺特}\Term{regular-scheme}{正则概形}$X$，环$\operatorname{Gr}'_{\mathrm{top}}(X)$是尚待建立的\Term{chow-ring}{Chow环}$A(X)$的一个方便的替代物。它相对于后者还有一个优点：[RRR]中某些计算直接在这个环内进行，并得到了重要公式；即使$X$是域上的\Term{smooth}{光滑}\Term{projective-scheme}{射影概形}，目前也还没有在$A(X)$本身中证明这些公式（第XIV讲4.3、4.4）。不过，要充分说明这一观点的合理性，仍须解决上面提到的$K(X)$的\Term{topological-filtration}{拓扑滤过}与其\Term{ring-structure}{环结构}的相容性问题。这个问题似乎本质上与Chow的\Term{moving-lemma}{移动引理}有关；而我们声称可以避开的\Term{chow-ring}{Chow环}理论，其关键技术成分恰恰就是这个\Term{lemma}{引理}。

\textbf{2.3.} 为避开此问题，在$K(X)$上引入第二个\Term{filtered}{滤过}，它比\Term{topological-filtration}{拓扑滤过}更细，并可由$K(X)$的\Term{augmented-lambda-ring}{增广$\lambda$环}结构以纯\Term{algebra}{代数}方式描述。其定义来自[RRR]中$x\in K(X)$的\Term{chern-class}{陈类}在$\operatorname{Gr}'_{\mathrm{top}}(X)$中的表达式：
\[c_i(x)=\gamma^i(x-\varepsilon(x))\pmod{\operatorname{Filt}^{i+1}_{\mathrm{top}}K(X)},\tag{2.3}\]
其中$\gamma^i$表示$\lambda^j$（$j\leq i$）的某个\Term{linear-combination}{线性组合}，上引文献已将它明确写出，$\varepsilon$为\Term{augmentation}{增广}。可以证明，$\gamma^i(x-\varepsilon(x))$的\Term{topological-filtration}{拓扑滤过}确实$\geq i$，因而(2.3)有意义。于是，用\Term{ideal}{理想}$\operatorname{Filt}^iK(X)$定义$K(X)$的\Term{lambda-filtration}{$\lambda$滤过}，其中$\operatorname{Filt}^iK(X)$由形如
\[\gamma^{i_1}(x_1-\varepsilon(x_1))\cdots\gamma^{i_k}(x_k-\varepsilon(x_k)),
\qquad i_1+\cdots+i_k\geq i\]
的表达式的和组成。此\Term{filtered-category}{滤过的}\Term{associated-graded-ring}{相伴分次环}记为$\operatorname{Gr}'(X)$，它将在本讨论班中代替$X$的\Term{chow-ring}{Chow环}。此外，可以证明，在这里考虑的情形中，\Term{canonical-homomorphism}{典范同态}
\[\operatorname{Gr}'(X)\longrightarrow\operatorname{Gr}'_{\mathrm{top}}(X)\tag{2.4}\]
\SourcePage{15}
模去挠以后是\Term{isomorphism}{同构}。这说明将该环张量以$\mathbb Q$以后，确实可以在各个方面起到\Term{chow-ring}{Chow环}张量以$\mathbb Q$的\Term{action}{作用}。另一方面，我们的构造恰好使公式(2.3)（其中以$\operatorname{Filt}^i$代替$\operatorname{Filt}^i_{\mathrm{top}}$）定义一个取值于$\operatorname{Gr}'(X)$的\Term{chern-class}{陈类}理论；再借助\Term{homomorphism}{同态}(2.4)，即得到$\operatorname{Gr}'_{\mathrm{top}}(X)$中的\Term{chern-class}{陈类}。利用$K(X)$是[RRR]意义下的特殊$\lambda$环（即第V讲意义下的$\lambda$环）这一事实（第VI讲），可以证明（第V讲）这个\Term{chern-class}{陈类}概念满足[RRR]所回顾的所有通常的\Term{form}{形式}性质。

\textbf{2.4.} $\operatorname{Gr}'(X)$的定义相对于$A(X)$与$\operatorname{Gr}'_{\mathrm{top}}(X)$，有一个很可贵的优点：它对完全任意的\Term{scheme}{概形}$X$，更\Term{generic}{一般}地，对任意\Term{locally-ringed-topos}{局部赋环拓扑斯}$X$仍然有意义；此时以$X$上的\Term{locally-free-module}{局部自由模}构造类群$K(X)$，它仍是\Term{augmentation}{增广}$\lambda$环，因而具有相应的\Term{filtered}{滤过}，可以定义\Term{associated-graded-ring}{相伴分次环}$\operatorname{Gr}'(X)$。但这种\Term{generic}{一般}性也有代价：如果不愿忽略挠现象，即不张量以$\mathbb Q$，则$\operatorname{Gr}'(X)$的性质在若干方面不如另两个候选者令人满意（第XIV讲第4节）。例如，再设$f:X\longrightarrow Y$为具有\Term{ample-invertible-module}{丰沛可逆模}的\Term{noetherian}{诺特}\Term{regular-scheme}{正则概形}之间的\Term{proper-morphism}{固有态射}，则\Term{generic}{一般}并不能断言，\Term{homomorphism}{同态}$f_*:K(X)\longrightarrow K(Y)$将$\operatorname{Filt}^iK(X)$映入$\operatorname{Filt}^{i-d}K(Y)$（其中$d$为$X$在$Y$上的\Term{relative-dimension}{相对维数}），从而也不能断言它在\Term{associated-graded-ring}{相伴分次环}上定义从$\operatorname{Gr}'(X)$到$\operatorname{Gr}'(Y)$的\Term{degree}{次数}$-d$\Term{homomorphism}{同态}。不过，张量以$\mathbb Q$后可以证明这一点（第VII讲），从而得到\Term{degree}{次数}$-d$的\Term{homomorphism}{同态}
\[f_*:\operatorname{Gr}'(X)_{\mathbb Q}\longrightarrow\operatorname{Gr}'(Y)_{\mathbb Q},\tag{2.5}\]
它起着“圈的\Term{direct-image}{直像}”\Term{homomorphism}{同态}$A(X)\longrightarrow A(Y)$的\Term{action}{作用}。

这样，当不假定存在底域，而$f:X\longrightarrow Y$是具有\Term{ample-invertible-module}{丰沛可逆模}的\Term{noetherian}{诺特}\Term{regular-scheme}{正则概形}之间的\Term{proper-morphism}{固有态射}时，我们就已具备赋予公式(1.2)意义的全部结构成分。
\SourcePage{16}
\SGAsection{3}{}
现在说明，去掉$X$、$Y$的\Term{regularity}{正则性}假设，而代以关于\Term{morphism}{态射}$f$的一个\Term{local-scheme}{局部}条件之后，怎样仍可赋予公式(2.2)意义；我们将明确说明这个条件。

\textbf{3.1.} 第一个困难在于，作为[RRR]中理论出发点的观察，即$X$上的层的类群$K(X)$，无论用\Term{coherent-sheaf}{凝聚层}还是用\Term{locally-free-module}{局部自由层}来定义都相同，本质上依赖于$X$的\Term{regularity}{正则性}假设。当$X$是\Term{generic}{一般}\Term{scheme}{概形}时（为简单起见，不妨设为\Term{noetherian-scheme}{诺特概形}），必须区分这两个群，分别记为$K_{\bullet}(X)$与$K'(X)$：前者关于$X$对\Term{proper-morphism}{固有态射}是\Term{covariant}{协变}的，后者关于$X$对任意\Term{morphism}{态射}都是\Term{contravariant}{反变}的。注意，$K'(X)$自然具有\Term{ring-structure}{环结构}（甚至是\Term{augmentation}{增广}$\lambda$\Term{ring-structure}{环结构}），而$K_{\bullet}(X)$\Term{generic}{一般}不再具有\Term{ring-structure}{环结构}，只具有$K'(X)$\Term{module-structure}{模结构}。从$K'(X)$到$K_{\bullet}(X)$的\Term{homomorphism}{同态}$x\longmapsto x\operatorname{cl}_{\bullet}(\mathcal O_X)$（其中$\operatorname{cl}_{\bullet}(F)$表示$F$在$K_{\bullet}(X)$中的类）\Term{generic}{一般}不再是\Term{isomorphism}{同构}。在广义黎曼--罗赫\Term{theorem}{定理}的表述中，我们只使用$K'(X)$、$K'(Y)$，不使用$K_{\bullet}(X)$、$K_{\bullet}(Y)$。如2.3所说，这些\Term{augmentation}{增广}$\lambda$环配上其自然的$\lambda$\Term{filtered}{滤过}，产生\Term{graded-ring}{分次环}$\operatorname{Gr}'(X)$、$\operatorname{Gr}'(Y)$以及“\Term{chern-class}{陈类}”\Term{homomorphism}{同态}
\[c_i:K'(X)\longrightarrow\operatorname{Gr}'(X),\]
$Y$的情形也一样。但困难出现在定义\Term{direct-image}{直像}\Term{homomorphism}{同态}
\[f_*:K'(X)\longrightarrow K'(Y)\tag{3.1}\]
之时。它不能再由朴素的公式
\[f_*(\operatorname{cl}'(F))=\sum_i(-1)^i\operatorname{cl}'(R^if_*(F))\]
来定义，因为即使$F$\Term{locally-free}{局部自由}，模$R^if_*(F)$\Term{generic}{一般}也不再\Term{locally-free}{局部自由}，甚至未必具有\Term{finite}{有限}\Term{tor-dimension}{Tor维数}，故上式失去了意义。

\Term{derived-category}{导出范畴}的观点为这一困难提供了解决方法。事实上，尽管$R^if_*(F)$\Term{generic}{一般}是“坏的”层，比方说具有无限\Term{tor-dimension}{Tor维数}，
\SourcePage{17}
因而无法为之赋予$K'(Y)$中的类，但$Y$上的\Term{complex}{复形}$\mathrm Rf_*(F)$，即$F$在\Term{derived-category}{导出范畴}意义下的\Term{direct-image}{直像}，却往往非常好；上述$R^if_*(F)$是它的\Term{cohomology-sheaf}{上同调层}。确切地说，设$F$在$\mathcal O_Y$上具有\Term{finite}{有限}\Term{tor-dimension}{Tor维数}；当$\mathcal O_X$在$\mathcal O_Y$上具有\Term{finite}{有限}\Term{tor-dimension}{Tor维数}且$F$\Term{locally-free}{局部自由}时，情况就是如此。于是可以证明（第III讲），\Term{complex}{复形}$\mathrm Rf_*(F)$是“\Term{perfect-complex}{完美的}”，这里是指：其\Term{cohomology-sheaf}{上同调层}\Term{coherent}{凝聚}（这一点原本就知道），而且其整体\Term{tor-dimension}{Tor维数}\Term{finite}{有限}；\Term{equivalence}{等价}地说，在$Y$的每个\Term{affine-open}{仿射开集}$U$上，这个\Term{complex}{复形}在\Term{derived-category}{导出范畴}$D(U)$中\Term{isomorphism}{同构}于一个\Term{degree}{次数}\Term{bounded}{有界}、每一次都是\Term{finite-type}{有限型}\Term{locally-free-module}{局部自由模}的\Term{complex}{复形}。从\Term{homological-algebra}{同调代数}的观点看，\Term{perfect-complex}{完美复形}与\Term{locally-free-module}{局部自由模}“同样好”，是后者的自然推广。它们还构成\Term{derived-category}{导出范畴}$D(Y)$的一个“\Term{triangulated-category}{三角子范畴}”，即对\Term{morphism}{态射}的锥（或“\Term{mapping-cylinder}{映射柱}”）\Term{stable}{稳定}的\Term{subcategory}{子范畴}。

另一方面，对每个\Term{triangulated-category}{三角范畴}$C$，都可赋予一个类群$K(C)$，它是\Term{abelian-category}{阿贝尔范畴}的类群构造的自然变体：取以$C$的\Term{object}{对象}为\Term{generator}{生成元}的\Term{free-group}{自由群}，再除去对于每个\Term{distinguished-triangle}{可区别三角形}
\[
\begin{tikzcd}[column sep=small,row sep=small]
& K'' \arrow[dl] & \\
K' \arrow[rr] && K \arrow[ul]
\end{tikzcd}
\]
所给出的关系
\[\operatorname{cl}(K)-\operatorname{cl}(K')-\operatorname{cl}(K'').\]
在当前情形下，代替[RRR]中基本观察的是如下事实：当\Term{scheme}{概形}$Y$具有\Term{ample-invertible-module}{丰沛可逆模}时，$Y$上每个\Term{perfect-complex}{完美复形}$L_{\bullet}$在$D(Y)$中整体地（而不只\Term{local-scheme}{局部}地）\Term{isomorphism}{同构}于一个\Term{degree}{次数}\Term{bounded}{有界}、每一次都是\Term{finite-type}{有限型}\Term{locally-free-module}{局部自由模}的\Term{complex}{复形}$L'_{\bullet}$（第II讲）。由此容易\Term{pushout}{推出}（第IV讲），\Term{canonical-homomorphism}{典范同态}
\[K'(X)\longrightarrow K(\operatorname{Parf}(Y)),\]
其中$\operatorname{Parf}(Y)$是$Y$上的\Term{perfect-complex}{完美复形}所成的\Term{triangulated-category}{三角范畴}，是一个\Term{isomorphism}{同构}。因而对$Y$上每个\Term{perfect-complex}{完美复形}$L_{\bullet}$，都可赋予一个类
\[\operatorname{cl}'(L_{\bullet})\in K'(Y),\]
\SourcePage{18}
它无非就是
\[\operatorname{cl}'(L_{\bullet})=\sum_i(-1)^i\operatorname{cl}'(L'_i),\]
其中$L'_{\bullet}$如上，故$L'_i$确实是$Y$上的\Term{locally-free-module}{局部自由模}，可以取其在$K'(Y)$中的类。

据此，只要$f:X\longrightarrow Y$是\Term{noetherian-scheme}{诺特概形}之间的\Term{proper-morphism}{固有}且\Term{finite}{有限}\Term{tor-dimension}{Tor维数}的\Term{morphism}{态射}，并且$Y$有\Term{ample-invertible-module}{丰沛可逆模}，就可定义\Term{homomorphism}{同态}(3.1)，只须对$X$上每个\Term{locally-free-module}{局部自由模}$F$置
\[f_*(\operatorname{cl}'(F))=\operatorname{cl}'(\mathrm Rf_*(F))\tag{3.2}\]
即可（此时$\mathrm Rf_*(F)$确实是$Y$上的\Term{perfect-complex}{完美复形}，右边有定义）。当$X$也有\Term{ample-invertible-module}{丰沛可逆模}时，$K'(X)$同样可以解释为$K(\operatorname{Parf}(X))$；\Term{homomorphism}{同态}(3.1)无非是由
\[\mathrm Rf_*:\operatorname{Parf}(X)\longrightarrow\operatorname{Parf}(Y)\]
所诱导的\Term{triangulated-category}{三角范畴}之间的\Term{exact-functor}{正合函子}，按\Term{functorial}{函子性}给出的\Term{homomorphism}{同态}。

\textbf{3.2.} 要赋予(1.2)意义，还须对$f$加上一个条件，使\Term{homomorphism}{同态}(2.5)，即
\[f_*:\operatorname{Gr}'(X)_{\mathbb Q}\longrightarrow\operatorname{Gr}'(Y)_{\mathbb Q},\tag{3.3}\]
有意义；亦即保证刚才定义的\Term{homomorphism}{同态}(3.1)，在容许\Term{translation}{平移}\Term{degree}{次数}并模去挠之后，与这些环的\Term{filtered}{滤过}相容。最后，还须定义一个“\Term{virtual-relative-tangent-bundle}{虚相对切丛}”元素
\[T_f\in K'(X).\]
对于最后一点，重新采用2.1提出的定义，可见应作一个保证\Term{conormal-module}{余法模}$N_{X/X'}$\Term{locally-free}{局部自由}的假设（沿用该处记号）。最自然的假设是$i$为“\Term{regular-immersion}{正则浸入}”，即它将$X$认作$X'$的一个\Term{subscheme}{子概形}，其\Term{ideal-of-definition}{定义理想}\Term{local-scheme}{局部}地可以由$\mathcal O_{X'}$的\Term{section}{截面}构成的$\mathcal O_{X'}$\Term{regular-sequence}{正则序列}来定义。这是一个纯\Term{local-nature}{局部性质}的假设，实际上与将$f$\Term{resolution}{分解}为\Term{immersion}{浸入}后接\Term{smooth-morphism}{光滑态射}的方式无关（第VIII讲）；如同前面，由公式(2.1)定义的类$T_f$也与这种选择无关。当刚才的假设满足时，称$f$是“\Term{lci-morphism}{局部完全交态射}”\footnote{我们也允许在此术语中省略“\Term{local-scheme}{局部}”一词，这不会造成混淆。}。这样的\Term{morphism}{态射}还具有\Term{locally-finite}{局部有限}\Term{tor-dimension}{Tor维数}。
\SourcePage{19}
若$f$是\Term{lci-morphism}{局部完全交态射}，则由公式
\[d(f)(x)=\dim_x\bigl(f'^{-1}f(x)\bigr)-\operatorname{codim}_x(X,X')\]
引入自然的“\Term{virtual-relative-dimension}{虚相对维数}”$d(f)$概念。它是关于点$x\in X$的\Term{locally-constant}{局部常值}整数值函数，故当$X$\Term{connected}{连通}时为常数；它也与\Term{resolution}{分解}(2.1)的选择无关。当这个\Term{virtual-relative-dimension}{虚相对维数}$d$为常数、$f$\Term{proper}{固有}且$X$、$Y$均有\Term{ample-invertible-module}{丰沛可逆模}时，可以证明（第VIII讲），尽管这绝非平凡的结果，仍有
\[f_*\bigl(\operatorname{Filt}^iK'(X)_{\mathbb Q}\bigr)
\subset\operatorname{Filt}^{i-d}K'(Y)_{\mathbb Q},\tag{3.4}\]
从而得到\Term{degree}{次数}为$-d$的\Term{homomorphism}{同态}(3.3)。

\textbf{3.3.} 因此，对于均具有\Term{ample-invertible-module}{丰沛可逆模}的\Term{noetherian-scheme}{诺特概形}$X$、$Y$之间的\Term{proper}{固有}\Term{lci-morphism}{局部完全交态射}，我们已经具备了赋予公式(1.2)意义的全部结构成分。稍加注意，这个公式事实上在不作任何\Term{noetherian}{诺特}性假设时仍有意义。

这个\Term{generic}{一般}\Term{form}{形式}的黎曼--罗赫公式将在第VIII讲证明。读者将看到，本讨论班给出的证明，结合了[RRR]所载的、引言中提到的两种证明的成分。必须采用第一种证明的方法（它比Borel--Serre所采用的第二种证明需要更多的\Term{form}{形式}理论），是因为需要证明包含关系(3.4)；在[RRR]和Borel--Serre [2]的较狭窄框架内，可以使用\Term{chow-ring}{Chow环}或$\operatorname{Gr}'_{\mathrm{top}}$环，因而并不存在这一问题。我们从第二种证明保留了这样的做法：当$f$是\Term{immersion}{浸入}时，引入$Y$沿$X$的\Term{blow-up}{爆破概形}。一旦第XIV讲第1节提出的、在\Term{derived-category}{导出范畴}$D^-(Y)$中引入\Term{exterior-power}{外幂}运算的问题得到解决，这个技术手段无疑就会消失，而\Term{immersion}{浸入}情形的证明也将推广到不再假定$Y$上存在\Term{ample-invertible-module}{丰沛可逆模}的情形。
\SGAsection{4}{}
最后还须略加说明：当不假定$X$、$Y$上存在\Term{ample-invertible-module}{丰沛可逆模}时，怎样仍可赋予(1.2)意义。
\SourcePage{20}
\textbf{4.1.} 首先注意，对于任意\Term{scheme}{概形}$X$（或更\Term{generic}{一般}地，任意\Term{local-scheme}{局部}\Term{ringed-topos}{赋环拓扑斯}），似乎毫无疑问，代替[RRR]中的环$K(X)$的“正确”\Term{invariants}{不变量}，应当是3.1所考察的、用$X$上\Term{perfect-complex}{完美复形}定义的群$K(\operatorname{Parf}(X))$，而非此前用$X$上\Term{locally-free-module}{局部自由模}定义的环$K'(X)$。现在将前者记作$K'(X)$，以便与3.1中的“朴素”群区分，后者记为$K'_{\text{naïf}}(X)$。对于\Term{proper}{固有}、\Term{finite}{有限}\Term{tor-dimension}{Tor维数}的\Term{morphism}{态射}$f:X\longrightarrow Y$，必须定义\Term{direct-image}{直像}\Term{homomorphism}{同态}(3.1)，这一需要迫使我们采取上述观点；这里不妨先设\Term{scheme}{概形}是\Term{noetherian}{诺特}的（非\Term{noetherian}{诺特}情形要求$f$满足略强的\Term{locally-finite}{局部有限}性性质，参见第III讲）。采用现在的$K'(X)$定义后，\Term{homomorphism}{同态}$f_*$确实可直接由公式(3.2)定义，其中$F$现在是$X$上的任意\Term{perfect-complex}{完美复形}，因此$\mathrm Rf_*(F)$也是\Term{perfect}{完美}的。

\textbf{4.2.} 然而，这个新的$K'(X)$定义带来一个新的\Term{homological-algebra}{同调代数}问题：如何在$K'(X)$上定义一个$\lambda$结构，使之与\Term{derived-category}{导出范畴}中的\Term{tensor-product}{张量积}所给出的自然\Term{ring-structure}{环结构}相容。（注意：若$L_{\bullet}$、$L'_{\bullet}$为两个\Term{perfect-complex}{完美复形}，则其全\Term{tensor-product}{张量积}$L_{\bullet}\otimes^{\mathrm L}L'_{\bullet}$仍为\Term{perfect-complex}{完美复形}。）更确切地说，应在\Term{derived-category}{导出范畴}$D^-(X)$中定义“\Term{exterior-power}{外幂}”运算，并使它将\Term{perfect-complex}{完美复形}变为\Term{perfect-complex}{完美复形}。目前尚未澄清这个问题；如引言所说，这构成本讨论班最严重的缺憾。一旦解决这个问题，仿照2.3，就得到$K'(X)$上的自然\Term{filtered}{滤过}，进而\Term{constructible}{可构造}\Term{associated-graded-ring}{相伴分次环}$\operatorname{Gr}'(X)$，使之起\Term{chow-ring}{Chow环}的\Term{action}{作用}，并定义\Term{chern-class}{陈类}理论
\[c_i:K'(X)\longrightarrow\operatorname{Gr}'(X).\]

\textbf{4.3.} 为赋予\Term{noetherian-scheme}{诺特概形}之间的\Term{proper}{固有}\Term{lci-morphism}{局部完全交态射}$f:X\longrightarrow Y$的公式(1.2)意义，还必须在该公式的陈述中包含非平凡的包含关系(3.4)（其中$d$为\Term{virtual-relative-dimension}{虚相对维数}），以便定义(3.3)；最后，还须定义虚相对切类
\SourcePage{21}
\[T_f\in K'(X).\]
当$f$有\Term{resolution}{分解}(2.1)，即先\Term{immersion}{浸入}一个相对\Term{smooth-scheme}{光滑概形}时，$T_f$的定义没有带来新的困难。此外，在此情形下，还可以采取更符合“\Term{derived-category}{导出范畴}”精神的表述：考虑众所周知的\Term{canonical-homomorphism}{典范同态}
\[N_{X/X'}\xrightarrow{d}\Omega^1_{X'/Y}\otimes_{\mathcal O_{X'}}\mathcal O_X,\tag{4.1}\]
它由外\Term{differential}{微分}
\[d_{X'/X}:J\subset\mathcal O_{X'}\longrightarrow\Omega^1_{X'/Y}\]
诱导，再考虑借助(4.1)“\Term{extension-by-zero}{以零延拓}”所定义的\Term{chain-complex}{链复形}
\[L_{\bullet}^{X/Y}\in\operatorname{Ob}D(X),\tag{4.2}\]
这个\Term{complex}{复形}在\Term{degree}{次数}1、0处分别等于\Term{homomorphism}{同态}(4.1)的源、靶。公式(2.2)于是化为更令人满意的\Term{form}{形式}
\[\check T_f=\operatorname{cl}'(L_{\bullet}^{X/Y}),\tag{4.3}\]
当$f$是\Term{lci-morphism}{局部完全交态射}时，这个公式在$K'(X)$中有意义，因为此时\Term{complex}{复形}$L_{\bullet}^{X/Y}$显然是\Term{perfect}{完美}的。关于$T_f$不依赖于$f$的所选\Term{resolution}{分解}的断言，现在可以容易地精确表述为\Term{complex}{复形}$L_{\bullet}^{X/Y}$的\Term{invariant}{不变}性断言：在\Term{derived-category}{导出范畴}$D(X)$中，模去\Term{canonical-isomorphism}{典范同构}，这个\Term{complex}{复形}也与所选\Term{resolution}{分解}无关（第VIII讲）。

\textbf{4.4.} 当$f:X\longrightarrow Y$是任意\Term{scheme-morphism}{概形态射}时，仍\Term{constructible}{可构造}\Term{derived-category}{导出范畴}$D(X)$中的一个\Term{canonical}{典范}\Term{object}{对象}$L_{\bullet}^{X/Y}$，它在唯一\Term{isomorphism}{同构}的意义下确定；在$X$的每个映到$Y$的某\Term{affine-open}{仿射开集}之上的\Term{affine-open}{仿射开集}上，它都\Term{isomorphism}{同构}于按上一段原则构造的\Term{complex}{复形}（但在\Term{generic}{一般}情形中，应将“\Term{smooth}{光滑}”换成“\Term{formally-smooth}{形式光滑}”，使得环$A$上的任意\Term{algebra}{代数}$B$都\Term{representable-functor}{可表示}成一个\Term{formally-smooth}{形式光滑}\Term{algebra}{代数}的商，例如\Term{polynomial-algebra}{多项式代数}的商）。这样的定义不能单靠\Term{affine}{仿射}构造的“\Term{gluing}{粘合}”来完成，因为\Term{derived-category}{导出范畴}$D(X)$的\Term{object}{对象}本质上不能\Term{gluing}{粘合}。
\SourcePage{22}
下面给出$L_{\bullet}^{X/Y}$的\Term{generic}{一般}构造，它更\Term{generic}{一般}地适用于任意（交换）\Term{ringed-topos}{赋环拓扑斯}的\Term{morphism}{态射}$f:X\longrightarrow Y$。置$B=\mathcal O_X$、$A=f^{-1}(\mathcal O_Y)$，于是$A$是\Term{space}{空间}$X$上的一个环，$B$是一个$A$\Term{algebra}{代数}。设$T\longrightarrow B$为\Term{sheaf}{集合层}的一个\Term{homomorphism}{同态}，它将$B$生成为$A$\Term{algebra}{代数}，即相应的$A$\Term{algebra-homomorphism}{代数同态}
\[\phi:C=A[T]\longrightarrow B\]
是\Term{set-sheaf}{集合层}的\Term{epimorphism}{满同态}；其中$A[T]$表示由$T$生成的自由交换$A$\Term{algebra}{代数}（即\Term{presheaf}{预层}$U\longmapsto A(U)[T(U)]$的\Term{associated-sheaf}{伴随层}，$A(U)[T(U)]$是由$T(U)$生成的自由交换$A(U)$\Term{algebra}{代数}）。设$J$为$\phi$的核，如4.3中那样，考虑由自然\Term{homomorphism}{同态}
\[J/J^2\xrightarrow{d}\Omega^1_{C/A}\otimes_C B\]
定义的\Term{length}{长度}为1的\Term{chain-complex}{链复形}。不难看出，在\Term{category}{范畴}$D(X)$中模去\Term{canonical-isomorphism}{典范同构}后，这样构造的\Term{complex}{复形}与$T$以及$T\longrightarrow B$的选择无关，因此确实应称为$X$在$Y$上的\Term{cotangent-complex}{相对余切复形}$L_{\bullet}^{X/Y}$。也可验证，在$X$的每个映到$Y$的某\Term{affine-open}{仿射开集}之上的\Term{affine-open}{仿射开集}上，它都由4.3的方法给出（模去\Term{canonical-isomorphism}{典范同构}）。

当$f$\Term{local-complete-intersection}{局部完全交}时，还可看出所得到的\Term{complex}{复形}甚至是\Term{perfect}{完美}的，因此仍可用公式(4.3)在$K'(X)$中定义元素$T_f$。须指出，在这里\Term{generic}{一般}也无法把$T_f$定义为$K'_{\text{naïf}}(X)$中的元素，也没有理由认为$T_f$在$D(X)$中整体地\Term{isomorphism}{同构}于一个各\Term{set-component}{分量}均为\Term{finite-type}{有限型}\Term{locally-free-module}{局部自由模}的\Term{bounded-complex}{有界复形}。因此，这又进一步说明4.1所主张的观点是合理的：应采用“精细”的$K'(X)$，而采用$K'_{\text{naïf}}(X)$不如前者适宜。

\textbf{4.5.} 撇开4.2提出的问题，我们因而再次具备了赋予任意\Term{scheme}{概形}之间的\Term{proper}{固有}\Term{lci-morphism}{局部完全交态射}的公式(1.2)意义的全部结构成分（事实上，前几节暗中采用的\Term{noetherian}{诺特}性假设，可以不遇到实质困难地消去）。遗憾的是，这并不意味着，即使暂且撇开上述问题，该公式就已得证；即使$Y$是\Term{characteristic}{特征}$p>0$的\Term{algebraically-closed-field}{代数闭域}的谱，
\SourcePage{23}
而$X$是不妨说三维的\Term{proper}{固有}\Term{smooth}{光滑}\Term{algebraic-scheme}{代数概形}，情况也是如此。参见第XIV讲第2节的评述。
\SGAsection{5}{}
按第4节所述的\Term{form}{形式}，黎曼--罗赫公式(1.2)对于\Term{complex-analytic-space}{复解析空间}（或Tate意义下的\Term{rigid-analytic-space}{刚性解析空间}）之间的\Term{proper}{固有}\Term{lci-morphism}{局部完全交态射}也有意义。当然，[RRR]给出的证明只有在额外假定\Term{morphism}{态射}$f$\Term{projective}{射影}时，才适用于这一情形；如同\Term{algebra}{代数}情形，处理\Term{generic}{一般}情形似乎需要一个本质上新的思想。这里所考虑的公式取值于“解析\Term{chow-ring}{Chow环}”$\operatorname{Gr}'(X)$，更确切地说，取值于$\operatorname{Gr}'(X)_{\mathbb Q}$；它依4.2说明的原则构造（该处遇到的困难，由于处于\Term{characteristic-zero}{特征零}而消失，参见第XIV讲第1节）。

当$X$、$Y$\Term{nonsingular}{非奇异}时，还可以如下给出黎曼--罗赫公式的另一个版本。先设$X$为任意\Term{complex-analytic-space}{复解析空间}，以$X^{\mathrm{top}}$表示同一个\Term{space}{空间}配上\Term{continuous}{连续}复值函数层。于是有\Term{ringed-space}{赋环空间}的\Term{canonical-morphism}{典范态射}
\[X^{\mathrm{top}}\longrightarrow X;\]
由于\Term{functor}{函子}$K'(X)$关于可变\Term{ringed-space}{赋环空间}$X$显然\Term{contravariant}{反变}，它定义一个\Term{ring-homomorphism}{环同态}
\[K'(X)\longrightarrow K'(X^{\mathrm{top}}).\tag{5.1}\]
为简单起见，设$X^{\mathrm{top}}$\Term{compact}{紧致}；那么，如同具有\Term{ample-invertible-module}{丰沛可逆模}的\Term{scheme}{概形}情形（3.1中已经提及），可以看出（第II讲），$X^{\mathrm{top}}$上每个\Term{perfect-complex}{完美复形}都整体地\Term{isomorphism}{同构}于一个\Term{bounded}{有界}、每一次\Term{locally-free}{局部自由}的\Term{complex}{复形}，即$X^{\mathrm{top}}$上复\Term{vector-bundle}{向量丛}的\Term{complex}{复形}。因此，$K'(X^{\mathrm{top}})$\Term{canonical-isomorphism}{典范同构}于\Term{topology}{拓扑}学家熟知的群$K(X^{\mathrm{top}})$ [1]，后者用\Term{compact}{紧致}\Term{space}{空间}$X^{\mathrm{top}}$上的复\Term{vector-bundle}{向量丛}定义。于是可以把\Term{homomorphism}{同态}(5.1)解释为从用解析结构定义的$K'(X)$到\Term{topology}{拓扑}学家熟知的环$K(X^{\mathrm{top}})$的\Term{homomorphism}{同态}。现在，若$f:X\longrightarrow Y$是\Term{compact}{紧致}\Term{nonsingular}{非奇异}\Term{analytic-space}{解析空间}之间的\Term{morphism}{态射}，由Atiyah--Hirzebruch [1]的工作，知道可为之赋予一个\Term{group-homomorphism}{群同态}
\SourcePage{24}
\[f_*^{\mathrm{top}}:K'(X^{\mathrm{top}})\longrightarrow K'(Y^{\mathrm{top}}).\tag{5.2}\]
同样，使用公式(3.2)所定义的(3.1)中的\Term{homomorphism}{同态}$f_*$（这里暗中使用了Grauert关于\Term{analytic-space}{解析空间}之间\Term{proper-morphism}{固有态射}的\Term{finiteness-theorem}{有限性定理}[4]），就得到一个自然\Term{homomorphism}{同态}组成的方块，其中只涉及$K'$型群（不涉及\Term{chow-ring}{Chow环}或\Term{cohomology-ring}{上同调环}这类环）：
\[
\begin{tikzcd}
K'(X) \arrow[r] \arrow[d,"f_*"'] & K'(X^{\mathrm{top}}) \arrow[d,"f_*^{\mathrm{top}}"] \\
K'(Y) \arrow[r] & K'(Y^{\mathrm{top}})
\end{tikzcd}
\tag{5.3}
\]
我们所说的黎曼--罗赫型公式，就是断言这个方块交换。须注意，这个断言与第4节所考虑的黎曼--罗赫公式有性质上的区别：前者不忽略挠现象，并且是在本质上\Term{discrete}{离散}的群$K'(Y^{\mathrm{top}})$中成立；后者则在某种“Chow”型群中成立，它不再具有\Term{discrete}{离散}性质，但相应地必须张量以$\mathbb Q$。注意，当$X$、$Y$是\Term{nonsingular}{非奇异}\Term{projective-variety}{射影簇}时，公式(5.3)无论如何已被证明[7]；第4节所考虑类型的公式也是如此，在此情形下已在[RRR]中证明。这是Atiyah--Singer指标\Term{theorem}{定理}（以及Shih [6]与Atiyah所给推广）的一个合理的变体，显然应有一个共同的推广。一旦证明(5.3)的\Term{commutativity}{交换性}，便可借助Atiyah--Hirzebruch [1]的“可微黎曼--罗赫\Term{theorem}{定理}”，立即\Term{pushout}{推出}复解析黎曼--罗赫公式的一个\Term{cohomology}{上同调}版本，即下列方块的\Term{commutativity}{交换性}：
\[
\begin{tikzcd}[column sep=large]
K'(X) \arrow[r,"{x\mapsto\operatorname{ch}(x)\operatorname{Todd}(T_X)}"] \arrow[d] & H^{2*}(X,\mathbb Q) \arrow[d] \\
K'(Y) \arrow[r,"{y\mapsto\operatorname{ch}(y)\operatorname{Todd}(T_Y)}"] & H^{2*}(Y,\mathbb Q).
\end{tikzcd}
\tag{5.4}
\]
这里的水平\Term{arrow}{箭头}来自\Term{cohomology}{上同调}\Term{chern-class}{陈类}；这些\Term{chern-class}{陈类}由\Term{homomorphism}{同态}(5.1)和通常的\Term{chern-class}{陈类}\Term{homomorphism}{同态}
\[K'(X^{\mathrm{top}})\longrightarrow H^{2*}(X,\mathbb Z)\]
得到。
\SourcePage{25}
第一条竖直\Term{arrow}{箭头}与(5.3)中的相同，第二条则是有理\Term{cohomology}{上同调}中的Gysin\Term{homomorphism}{同态}。
\SGAsection{6}{}
在以上讨论中，几乎只谈到了\Term{contravariant}{反变}\Term{invariants}{不变量}$K'(X)$，而未谈及3.1已指出其存在的\Term{covariant}{协变}\Term{invariants}{不变量}$K_{\bullet}(X)$。如果把$K'(X)$看作\Term{topological-space}{拓扑空间}整系数\Term{cohomology-ring}{上同调环}的类似物，那么就应把$K_{\bullet}(X)$看作整系数\Term{homology}{同调}的类似物；$K_{\bullet}(X)$的$K'(X)$\Term{module-structure}{模结构}于是对应于帽积运算\footnote{这个类比并非纯粹是措辞上的。在转到\Term{associated-graded}{相伴分次对象}$\operatorname{Gr}'(X)$、$\operatorname{Gr}_{\bullet}(X)$以后，可以用$\operatorname{Gr}_{\bullet}(X)$上的\Term{module-structure}{模结构}与“帽积”运算之间的一个明确的相容关系，将它精确化。}。当$X$\Term{regular}{正则}时，自然\Term{homomorphism}{同态}$K'(X)\longrightarrow K_{\bullet}(X)$为\Term{isomorphism}{同构}，这一事实应看作\Term{orientation}{定向}\Term{topological-manifold}{拓扑流形}上庞加莱\Term{duality-theorem}{对偶定理}的类似物：该\Term{theorem}{定理}断言，与基本\Term{homology}{同调}类作帽积给出从\Term{cohomology}{上同调}到\Term{homology}{同调}的\Term{isomorphism}{同构}。这些类比还提示，例如对\Term{finite}{有限}多面体$X$，也适宜定义一个\Term{covariant}{协变}性质的\Term{invariants}{不变量}$K_{\bullet}(X)$，使它与整系数\Term{homology}{同调}的关系（\Term{spectral-sequence}{谱序列}、\Term{chern-homomorphism}{陈同态}等等），类似于$K'(X)$与整系数\Term{cohomology}{上同调}之间的关系；似乎也不排除，可以用这些群$K_{\bullet}(X)$，给出可微黎曼--罗赫\Term{theorem}{定理}的某些变体，使之对比\Term{compact}{紧致}\Term{differentiable-manifold}{可微流形}更\Term{generic}{一般}的\Term{space}{空间}成立。

回到$X$为任意\Term{noetherian-scheme}{诺特概形}的情形，应在$K_{\bullet}(X)$上赋予一个\Term{increasing-filtration}{递增滤过}，令$\operatorname{Filt}_iK_{\bullet}(X)$由所有$\operatorname{cl}_{\bullet}(F)$生成，其中$F$是$X$上的\Term{coherent-sheaf}{凝聚层}，且$\dim\operatorname{supp}F\leq i$。我们将证明（第X讲），这个\Term{filtered}{滤过}与\Term{module-structure}{模结构}和$K'(X)$的\Term{filtered}{滤过}相容，即有公式
\[\operatorname{Filt}^iK'(X)\cdot\operatorname{Filt}_jK_{\bullet}(X)
\subset\operatorname{Filt}_{j-i}K_{\bullet}(X).\tag{6.1}\]
于是可以在\Term{filtered}{滤过}模$K_{\bullet}(X)$的\Term{associated-graded}{相伴分次对象}（记作$\operatorname{Gr}_{\bullet}(X)$）上赋予$\operatorname{Gr}'(X)$\Term{module-structure}{模结构}。其\Term{homogeneous-element}{齐次元素}还可以解释为$X$上的\Term{algebraic-cycle}{代数圈}的类（当$X$在底域$k$上\Term{finite-type}{有限型}时，相应的\Term{equivalence-relation}{等价关系}比\Term{rational-equivalence}{有理等价}更粗）。如果
\SourcePage{26}
\[f:X\longrightarrow Y\]
是\Term{noetherian-scheme}{诺特概形}之间的\Term{proper-morphism}{固有态射}，那么
\[f_*:K_{\bullet}(X)\longrightarrow K_{\bullet}(Y)\tag{6.2}\]
与\Term{filtered}{滤过}相容，并定义\Term{graded-group}{分次群}\Term{homomorphism}{同态}
\[f_*:\operatorname{Gr}_{\bullet}(X)\longrightarrow\operatorname{Gr}_{\bullet}(Y),\tag{6.3}\]
它满足一个“\Term{projection-formula}{投影公式}”，即它是$\operatorname{Gr}'(Y)$\Term{linear}{线性}的。此公式由\Term{homomorphism}{同态}(6.2)（它是$K'(Y)$\Term{linear}{线性}的）的相应\Term{projection-formula}{投影公式}，转到\Term{associated-graded}{相伴分次对象}而得。

按照本讨论班提出的观点，在未必\Term{regular}{正则}的\Term{noetherian-scheme}{诺特概形}$X$上，一个便于运用的\Term{intersection-theory}{相交理论}，应当结合研究\Term{contravariant}{反变}\Term{invariants}{不变量}$K'(X)$、$\operatorname{Gr}'(X)$与\Term{covariant}{协变}\Term{invariants}{不变量}$K_{\bullet}(X)$、$\operatorname{Gr}_{\bullet}(X)$；它们承担着本质不同的\Term{action}{作用}，而其性质彼此补充。例如，对某些非\Term{proper}{固有}\Term{fibration}{纤维化}，\Term{covariant}{协变}群$K_{\bullet}(X)$、$\operatorname{Gr}_{\bullet}(X)$更便于计算，因为可以为它们建立“\Term{homotopy}{同伦}\Term{exact-sequence}{正合序列}”（[RRR]与第IX讲）。例如，可证明（第IX讲）如下\Term{canonical-isomorphism}{典范同构}：
\[K_{\bullet}(X[t])\simeq K_{\bullet}(X),\qquad
\operatorname{Gr}_{\bullet}(X[t])\simeq\operatorname{Gr}_{\bullet}(X),\]
而当$X$非\Term{regular}{正则}时，对$K'$\Term{generic}{一般}不成立。另一方面，$K'(X)$与$\operatorname{Gr}'(X)$由于具有\Term{ring-structure}{环结构}，也便于计算。

当考虑域上的\Term{proper-scheme}{固有概形}$X$时，\Term{projection}{投影}
\[f^X:X\longrightarrow(\text{点})\]
定义一个“欧拉--庞加莱\Term{euler-poincare-characteristic}{示性数}”\Term{homomorphism}{同态}
\[\chi_X=f^X_*:K_{\bullet}(X)\longrightarrow K_{\bullet}(\text{点})=\mathbb Z,\tag{6.4}\]
它在\Term{subgroup}{子群}$\operatorname{Gr}_0(X)=\operatorname{Filt}_0K_{\bullet}(X)$上诱\Term{derive}{导出}“零圈的\Term{degree}{次数}”\Term{homomorphism}{同态}
\[\deg_X:\operatorname{Gr}_0(X)\longrightarrow\mathbb Z.\tag{6.5}\]
\Term{multiplication}{乘法}\Term{homomorphism}{同态}
\[\operatorname{Gr}^i(X)\times\operatorname{Gr}_i(X)\longrightarrow\operatorname{Gr}_0(X)\]
与\Term{degree}{次数}\Term{homomorphism}{同态}(6.5)复合，便给出一个\Term{pairing}{配对}
\[\operatorname{Gr}^i(X)\times\operatorname{Gr}_i(X)\longrightarrow\mathbb Z,\tag{6.6}\]
这说明在某些方面$\operatorname{Gr}'(X)$与$\operatorname{Gr}_{\bullet}(X)$承担着\Term{duality}{对偶}的\Term{action}{作用}，
\SourcePage{27}
正如\Term{cohomology}{上同调}与\Term{homology}{同调}一样。例如，若$f:X\longrightarrow Y$是\Term{proper}{固有}\Term{algebraic-scheme}{代数概形}之间的\Term{morphism}{态射}，则两个\Term{homomorphism}{同态}
\[f^*:\operatorname{Gr}'(Y)\longrightarrow\operatorname{Gr}'(X),\qquad
f_*:\operatorname{Gr}_{\bullet}(X)\longrightarrow\operatorname{Gr}_{\bullet}(Y)\]
在\Term{pairing}{配对}(6.6)的意义下互为\Term{form}{形式}\Term{transpose}{转置}。

\Term{homomorphism}{同态}(6.4)、(6.5)及由此得到的\Term{pairing}{配对}(6.6)，可以看作给定域$k$上\Term{proper}{固有}\Term{algebraic-scheme}{代数概形}的\Term{intersection-theory}{相交理论}中的数值关系的来源。第X讲将较详细地展开这个观点；它基本独立于以证明黎曼--罗赫\Term{theorem}{定理}为\Term{central}{中心的}第VI、VII、VIII讲，也独立于第IX讲。那里还会给出关于$K'(X)$、$\operatorname{Gr}'(X)$在\Term{specialization}{特殊化}下的行为的一些概念。
\SGAsection{7}{}
如同任何\Term{intersection-theory}{相交理论}一样，必须说明\Term{picard-group}{Picard群}（或按经典术语称为\Term{divisor-class-group}{除子类群}）在理论中的地位。在一个轻微\Term{restriction}{限制}下，第X讲将建立一个\Term{canonical-isomorphism}{典范同构}
\[c_1:\operatorname{Pic}(X)\xrightarrow{\sim}\operatorname{Gr}^1(X),\tag{7.1}\]
例如当$X$\Term{quasi-compact}{拟紧}致时，该\Term{restriction}{限制}自动满足。这一\Term{isomorphism}{同构}表明\Term{picard-group}{Picard群}在我们提出的理论中的确切\Term{action}{作用}。因此，第XII、XIII讲将更详细地研究\Term{picard-group}{Picard群}；它们在逻辑上独立于本讨论班正文。Kleiman在第XIII讲中研究\Term{numerical-equivalence}{数值等价}以及\Term{picard-group}{Picard群}的若干\Term{finiteness-theorem}{有限性定理}（这些\Term{theorem}{定理}曾在[6]中宣布而未给证明）。他使用由他本人及Mumford提出的纯数值性质的技术（不使用Picard\Term{functor}{函子}及其\Term{representability}{可表性}理论），明显比Grothendieck最初的证明简单。只有一个结果无法由Kleiman的方法处理（他又不得不使用它），即[6，C-07 (i)]；第XII讲专门给出这个结果的证明及其若干改进（包括Picard\Term{functor}{函子}的某些\Term{representability}{可表性}\Term{theorem}{定理}）。这最后一讲的思想和技术完全不同于本讨论班其余部分，而与[5]相联系；它在逻辑上独立于其余各讲。
\SourcePage{28}
最后，第XI讲的内容远比第XII、XIII讲\Term{generic}{一般}，并更接近本讨论班的总体精神；它在任意\Term{local-scheme}{局部}\Term{ringed-topos}{赋环拓扑斯}上研究一个值得注意的\Term{functor}{函子}，称为“\Term{determinant-functor}{行列式函子}”：
\[\det{}^*:\operatorname{Parf}_{\mathrm{is}}(X)\longrightarrow\operatorname{Inv}(X),\tag{7.2}\]
从$X$上\Term{perfect-complex}{完美复形}的\Term{category}{范畴}（\Term{morphism}{态射}只取\Term{derived-category}{导出范畴}$D(X)$意义下的\Term{isomorphism}{同构}），映到$X$上\Term{invertible-module}{可逆模}的\Term{category}{范畴}。这个\Term{functor}{函子}基本由如下要求描述：它应当具有“\Term{local-nature}{局部性质}”，并且对\Term{bounded}{有界}、各\Term{set-component}{分量}\Term{locally-free}{局部自由}的\Term{complex}{复形}$L_{\bullet}$，由公式
\[\det{}^*(L_{\bullet})\simeq\bigotimes_i\det(L_i)^{(-1)^i}\tag{7.3}\]
给出，其中$\det(L_i)$表示\Term{locally-free-module}{局部自由模}$L_i$的最高次\Term{exterior-power}{外幂}。虽然除第I讲外，这一讲在逻辑上同样独立于本讨论班其余部分，但它仍是本讨论班所提出、并主要在第I至XI讲中展开的总体“瑜伽”（yoga）的重要组成部分。
\SourcePage{29}
\SGAsection{}{参考文献}
\begin{description}
\item[{[1]}] M. Atiyah与F. Hirzebruch，\textit{Vector bundles and homogeneous spaces}（\Term{vector-bundle}{向量丛}与\Term{homogeneous-space}{齐性空间}），\textit{Symposia in Pure Mathematics}，第3卷，\textit{Differential Geometry}，第7--38页，1961年。
\item[{[2]}] A. Borel与J. P. Serre，\textit{Le théorème de Riemann--Roch}（黎曼--罗赫\Term{theorem}{定理}），\textit{Bull. Soc. math. France}，第86卷，第97--136页，1958年。
\item[{[3]}] H. Grauert，\textit{Ein Theorem der analytischen Garbentheorie}（解析层论中的一个\Term{theorem}{定理}），\textit{Pub. Math.}，第5号，第5--64页，1960年。
\item[{[4]}] A. Grothendieck，\textit{Classes de faisceaux et théorème de Riemann--Roch}（层的类与黎曼--罗赫\Term{theorem}{定理}），油印讲义，普林斯顿，1957年（引用时记作[RRR]，并重印于本讨论班第0讲的附录）。
\item[{[5]}] A. Grothendieck，\textit{Technique de descente et théorèmes d'existence en Géométrie Algébrique}（\Term{algebraic-geometry}{代数几何}中的\Term{descent}{下降}技术与存在性\Term{theorem}{定理}），收入\textit{Fondements de la Géométrie Algébrique}（\Term{algebraic-geometry}{代数几何}基础），1957--1962年Bourbaki讨论班选辑，Secrétariat Mathématique，巴黎。
\item[{[6]}] W. Shih，IHES\Term{topology}{拓扑}学讨论班，1966/67年。
\item[{[7]}] M. Atiyah与F. Hirzebruch，\textit{The Riemann--Roch theorem for analytic embeddings}（解析\Term{embedding}{嵌入}的黎曼--罗赫\Term{theorem}{定理}），\textit{Topology}，第1卷，第151--166页，1962年。
\end{description}
